Verified external inputs used below: nonintegral surgeries on knots in \(S^3\) are irreducible; surgery on a knot in \(S^3\) has at most three prime summands, and in the three-summand case two are lens spaces and one an integer homology sphere; hyperbolic L-space knots have reducing slope \(2g(K)-1\); torus and satellite knots satisfy the cabling conjecture; the standard cable-surgery formula; and the \(V_k\)/\(d\)-invariant formulas used in the correction-term identity. ([cambridge.org](https://www.cambridge.org/core/journals/mathematical-proceedings-of-the-cambridge-philosophical-society/article/abs/only-integral-dehn-surgeries-can-yield-reducible-manifolds/6453897351AFACFE2991721E7F89AB13))

```latex
\documentclass[11pt]{article}
\usepackage[T1]{fontenc}
\usepackage[margin=1in]{geometry}
\usepackage{amsmath,amsthm,amssymb,mathtools}
\usepackage{enumitem}

\newtheorem{theorem}{Theorem}
\newtheorem{lemma}[theorem]{Lemma}
\newtheorem{proposition}[theorem]{Proposition}
\newtheorem{remark}[theorem]{Remark}

\newcommand{\Spinc}{\operatorname{Spin}^{c}}
\newcommand{\HF}{HF}

\title{Dehn surgery on an \(L\)-space knot has at most two prime summands}
\author{}
\date{}

\begin{document}
\maketitle

\begin{abstract}
We prove that if \(K\subset S^{3}\) is an \(L\)-space knot and \(r\in \mathbb{Q}\), then
\(S^{3}_{r}(K)\) has at most two prime summands.
The only nontrivial case is a hypothetical reducible surgery on a hyperbolic
\(L\)-space knot.  By a theorem of Hom--Lidman--Zufelt, such a surgery must occur
at slope \(2g(K)-1\).  Assuming there are three prime summands, a theorem of Howie
reduces the situation to a connected sum \(Y_{p}\# Y_{q}\) with
\(|H_{1}(Y_{p})|=p\), \(|H_{1}(Y_{q})|=q\), \(1<p<q\), and \(pq=2g(K)-1\).
A correction-term identity coming from the Ni--Wu surgery formula then contradicts
the basic properties of the \(V_{s}\)-sequence of a nontrivial \(L\)-space knot.
\end{abstract}

\section{Statement and standard inputs}

We write \(S^{3}_{r}(K)\) for \(r\)-surgery on \(K\subset S^{3}\), and
\(g(K)\) for the Seifert genus.

\begin{theorem}\label{thm:main}
Let \(K\subset S^{3}\) be an \(L\)-space knot, and let \(r\in \mathbb{Q}\).
Then the surgered manifold \(S^{3}_{r}(K)\) has at most two prime summands.
\end{theorem}

We shall use the following standard results.

\begin{proposition}\label{prop:standard}
The following facts hold.
\begin{enumerate}[label=\textup{(\roman*)}]
\item
If \(r\notin \mathbb{Z}\), then \(S^{3}_{r}(J)\) is irreducible for every knot
\(J\subset S^{3}\) \emph{[GL87]}.

\item
If a manifold obtained by Dehn surgery on a knot in \(S^{3}\) has more than two
prime summands, then it has exactly three, and is of the form
\[
L(p,a)\# L(q,b)\# Y,
\]
where \(p,q>1\), \((p,q)=1\), and \(Y\) is an integer homology sphere
\emph{[How02, Corollary~5.3]}.

\item
If \(K\) is a hyperbolic positive \(L\)-space knot and \(r\) is a reducing slope,
then \(r=2g(K)-1\) \emph{[HLZ15, Theorem~1.3]}.

\item
Torus knots and satellite knots satisfy the cabling conjecture:
if \(J\subset S^{3}\) is torus or satellite and \(S^{3}_{r}(J)\) is reducible,
then \(J\) is a cable knot and \(r\) is its cabling slope
\emph{[Mos71, Sch90]}.

\item
If \(C_{u,v}(J)\) denotes the \((u,v)\)-cable of a knot \(J\), with
\(u\ge 2\) and \((u,v)=1\), then
\[
S^{3}_{uv}\!\bigl(C_{u,v}(J)\bigr)\cong L(u,v)\# S^{3}_{v/u}(J)
\]
\emph{[Gor83]}.
\end{enumerate}
\end{proposition}

\begin{proof}
These are precisely the cited results.
\end{proof}

\section{A correction-term identity for reducible surgeries}

For a knot \(K\), let \(V_{s}=V_{s}(K)\) denote the standard sequence appearing in the
Ni--Wu correction-term formula.
We shall use the following basic properties:
\begin{equation}\label{eq:V-basic}
V_{s}\ge V_{s+1}\quad\text{for all }s,\qquad
V_{s}=0\quad\text{for all }s\ge g(K),
\end{equation}
proved in \emph{[NW15, Lemma~2.4]}.
If \(K\) is a nontrivial \(L\)-space knot, then
\begin{equation}\label{eq:V-positive}
V_{s}>0\qquad (0\le s\le g(K)-1),
\end{equation}
see \emph{[HLZ15, proof of Theorem~1.3]}.

The next lemma is the key identity.  It is essentially the same computation as in
\emph{[BD22, Theorem~3.1]}, but we include the short proof for completeness.

\begin{lemma}\label{lem:identity}
Let \(K\subset S^{3}\), let \(p,q>1\) be coprime, and suppose
\[
S^{3}_{pq}(K)\cong Y_{p}\# Y_{q},
\]
where \(H_{1}(Y_{p})\cong \mathbb{Z}/p\mathbb{Z}\) and
\(H_{1}(Y_{q})\cong \mathbb{Z}/q\mathbb{Z}\).
Set
\[
M=\frac{(p-1)(q-1)}{2},
\qquad
\alpha(j)=\min\{j,pq-j\}.
\]
Then for every integer \(\ell\) with \(0\le \ell\le M\),
\[
\sum_{i=0}^{q-1}\Bigl(V_{\ell+i}-V_{\alpha(\ell+i+p)}\Bigr)=M-\ell.
\]
\end{lemma}

\begin{proof}
Fix \(0\le \ell\le M\).
Since
\[
0\le \ell+i \le M+q-1 < pq,
\qquad
0\le \ell+i+p \le M+q-1+p < pq,
\]
all indices below may be viewed as integers in \([0,pq-1]\).

Identify \(\Spinc\!\left(S^{3}_{pq}(K)\right)\cong \mathbb{Z}/pq\mathbb{Z}\).
Let
\[
\pi_{p}:\Spinc\!\left(S^{3}_{pq}(K)\right)\to \Spinc(Y_{p}),
\qquad
\pi_{q}:\Spinc\!\left(S^{3}_{pq}(K)\right)\to \Spinc(Y_{q})
\]
denote the projections.
Since \(|H_{1}(Y_{p})|=p\), we have \(\pi_{p}([t+p])=\pi_{p}([t])\) for every \(t\).
By additivity of correction terms under connected sum,
\begin{equation}\label{eq:additivity-step}
d\!\left(S^{3}_{pq}(K),[t+p]\right)-d\!\left(S^{3}_{pq}(K),[t]\right)
=
d\!\left(Y_{q},\pi_{q}([t+p])\right)-d\!\left(Y_{q},\pi_{q}([t])\right).
\end{equation}

Now sum \eqref{eq:additivity-step} over \(t=\ell,\ell+1,\dots,\ell+q-1\).
Because \(\pi_{q}\) restricts to a bijection from each of the sets
\[
\{[\ell],\dots,[\ell+q-1]\},
\qquad
\{[\ell+p],\dots,[\ell+p+q-1]\}
\]
onto \(\Spinc(Y_{q})\), the right-hand side sums to \(0\).  Hence
\begin{equation}\label{eq:sumzero}
\sum_{i=0}^{q-1}
\Bigl(
d(S^{3}_{pq}(K),[\ell+i+p])-d(S^{3}_{pq}(K),[\ell+i])
\Bigr)=0.
\end{equation}

For integer surgeries, the Ni--Wu formula gives
\[
d(S^{3}_{pq}(K),[s])=d(L(pq,1),[s])-2V_{\alpha(s)}
\qquad (0\le s<pq),
\]
see \emph{[BD22, Theorem~2.1]}, which is the integer specialization of
\emph{[NW15, Proposition~1.6]}.
Also, by \emph{[OS03, Proposition~4.8]},
\[
d(L(n,1),[u])=\frac{(2u-n)^{2}-n}{4n}\qquad (0\le u<n).
\]
Therefore
\[
d(L(pq,1),[\ell+i+p])-d(L(pq,1),[\ell+i])
=
\frac{2(\ell+i)+p-pq}{q}.
\]
Substituting into \eqref{eq:sumzero}, we obtain
\[
0
=
\sum_{i=0}^{q-1}
\left(
\frac{2(\ell+i)+p-pq}{q}
+
2V_{\ell+i}
-
2V_{\alpha(\ell+i+p)}
\right).
\]
Now
\[
\sum_{i=0}^{q-1}\frac{2(\ell+i)+p-pq}{q}
=
2\ell+q-1+p-pq
=
-2\left(\frac{(p-1)(q-1)}{2}-\ell\right)
=
-2(M-\ell).
\]
Thus
\[
0=-2(M-\ell)+2\sum_{i=0}^{q-1}\Bigl(V_{\ell+i}-V_{\alpha(\ell+i+p)}\Bigr),
\]
and the lemma follows.
\end{proof}

\section{The key obstruction and the proof of the main theorem}

\begin{proposition}\label{prop:key}
Let \(K\subset S^{3}\) be a nontrivial positive \(L\)-space knot of genus \(g\).
Suppose \(2g-1=pq\) with \(1<p<q\) and \((p,q)=1\).
Then there is no decomposition
\[
S^{3}_{2g-1}(K)\cong Y_{p}\# Y_{q}
\]
with \(|H_{1}(Y_{p})|=p\) and \(|H_{1}(Y_{q})|=q\).
\end{proposition}

\begin{proof}
Set
\[
M=\frac{(p-1)(q-1)}{2},
\qquad
d=\frac{q-p}{2}.
\]
Since \(2g-1=pq\) is odd, both \(p\) and \(q\) are odd, so \(d\in \mathbb{Z}_{>0}\).
Also
\[
g=\frac{pq+1}{2}=M+\frac{p+q}{2}=M+p+d.
\]

Apply Lemma~\ref{lem:identity} with \(\ell=M\).  We obtain
\begin{equation}\label{eq:ellM}
0=\sum_{i=0}^{q-1}\Bigl(V_{M+i}-V_{\alpha(M+i+p)}\Bigr).
\end{equation}

We now compute the indices \(\alpha(M+i+p)\).

If \(0\le i\le d-1\), then
\[
M+i+p\le M+d-1+p=g-1,
\]
hence
\[
\alpha(M+i+p)=M+i+p.
\]

If \(d\le i\le q-1\), then \(M+i+p\ge g\), so
\[
\alpha(M+i+p)=pq-(M+i+p)=M+q-1-i.
\]
Therefore \eqref{eq:ellM} becomes
\begin{equation}\label{eq:split}
0=
\sum_{i=0}^{d-1}\bigl(V_{M+i}-V_{M+p+i}\bigr)
+
\sum_{i=d}^{q-1}\bigl(V_{M+i}-V_{M+q-1-i}\bigr).
\end{equation}

For \(d\le i\le q-1-d\), the \(i\)-th term in the second sum of \eqref{eq:split}
cancels the \((q-1-i)\)-st term, because
\[
V_{M+i}-V_{M+q-1-i}= -\bigl(V_{M+q-1-i}-V_{M+i}\bigr).
\]
Hence \eqref{eq:split} reduces to
\[
0=\sum_{i=0}^{d-1}\bigl(V_{M+q-1-i}-V_{M+p+i}\bigr).
\]
Now
\[
M+q-1-i-(M+p+i)=q-p-1-2i=2d-1-2i>0
\qquad (0\le i\le d-1),
\]
so by the monotonicity \eqref{eq:V-basic} every summand is \(\le 0\).
But the summand with \(i=d-1\) is
\[
V_{M+q-d}-V_{M+p+d-1}=V_{g}-V_{g-1},
\]
because \(M+q-d=g\) and \(M+p+d-1=g-1\).
By \eqref{eq:V-basic} and \eqref{eq:V-positive},
\[
V_{g}=0,\qquad V_{g-1}>0,
\]
so this summand is strictly negative.
Hence the whole sum is strictly negative, a contradiction.
\end{proof}

\begin{proof}[Proof of Theorem~\ref{thm:main}]
If \(K\) is the unknot, then every surgery on \(K\) is a lens space or \(S^{1}\times S^{2}\),
so the conclusion is immediate.  Hence assume \(K\) is nontrivial.

If \(S^{3}_{r}(K)\) is irreducible, then it has exactly one prime summand.
So suppose \(S^{3}_{r}(K)\) is reducible.

Replacing \(K\) by its mirror if necessary, we may assume that \(K\) is a
\emph{positive} \(L\)-space knot; this does not change the number of prime summands,
since
\[
S^{3}_{r}(\overline{K})\cong -\,S^{3}_{-r}(K).
\]

\smallskip
\noindent\textbf{Case 1: \(K\) is torus or satellite.}
By Proposition~\ref{prop:standard}(iv), \(K\) is a cable knot and \(r\) is the cabling slope.
Write \(K=C_{u,v}(J)\) with \(u\ge 2\).
By Proposition~\ref{prop:standard}(v),
\[
S^{3}_{r}(K)=S^{3}_{uv}(C_{u,v}(J))
\cong L(u,v)\# S^{3}_{v/u}(J).
\]
Since \(u\ge 2\), the slope \(v/u\) is nonintegral; therefore
\(S^{3}_{v/u}(J)\) is irreducible by Proposition~\ref{prop:standard}(i).
Hence \(S^{3}_{r}(K)\) is a connected sum of at most two prime manifolds.

\smallskip
\noindent\textbf{Case 2: \(K\) is hyperbolic.}
By Proposition~\ref{prop:standard}(iii),
\[
r=2g(K)-1.
\]
Assume for contradiction that \(S^{3}_{r}(K)\) has more than two prime summands.
Then Proposition~\ref{prop:standard}(ii) gives
\[
S^{3}_{r}(K)\cong L(p,a)\# L(q,b)\# Y,
\]
where \(p,q>1\), \((p,q)=1\), and \(Y\) is an integer homology sphere.
After interchanging the two lens-space summands if necessary, assume \(1<p<q\).

Since
\[
H_{1}\!\left(S^{3}_{r}(K)\right)\cong \mathbb{Z}/r\mathbb{Z},
\]
we must have \(r=pq\).  Thus
\[
2g(K)-1=pq.
\]
Now group the last two summands together:
\[
S^{3}_{2g(K)-1}(K)\cong L(p,a)\#\bigl(L(q,b)\# Y\bigr).
\]
The first factor has first homology of order \(p\), and the second of order \(q\).
This contradicts Proposition~\ref{prop:key}.

The contradiction shows that the hyperbolic case also has at most two prime summands.
This completes the proof.
\end{proof}

\begin{thebibliography}{99}

\bibitem{BD22}
H.~Bodish and R.~DeYeso III,
\emph{Obstructing reducible surgeries: slice genus and thickness bounds},
arXiv:2209.01672.

\bibitem{GL87}
C.~McA. Gordon and J.~Luecke,
\emph{Only integral Dehn surgeries can yield reducible manifolds},
Math. Proc. Cambridge Philos. Soc. \textbf{102} (1987), no.~1, 97--101.

\bibitem{Gor83}
C.~McA. Gordon,
\emph{Dehn surgery and satellite knots},
Trans. Amer. Math. Soc. \textbf{275} (1983), no.~2, 687--708.

\bibitem{HLZ15}
J.~Hom, T.~Lidman, and N.~Zufelt,
\emph{Reducible surgeries and Heegaard Floer homology},
Math. Res. Lett. \textbf{22} (2015), no.~3, 763--788.

\bibitem{How02}
J.~Howie,
\emph{A proof of the Scott--Wiegold conjecture on free products of cyclic groups},
J. Pure Appl. Algebra \textbf{173} (2002), no.~2, 167--176.

\bibitem{Mos71}
L.~Moser,
\emph{Elementary surgery along a torus knot},
Pacific J. Math. \textbf{38} (1971), 737--745.

\bibitem{NW15}
Y.~Ni and Z.~Wu,
\emph{Cosmetic surgeries on knots in \(S^{3}\)},
J. Reine Angew. Math. \textbf{706} (2015), 1--17.

\bibitem{OS03}
P.~Ozsv\'ath and Z.~Szab\'o,
\emph{Absolutely graded Floer homologies and intersection forms for four-manifolds with boundary},
Adv. Math. \textbf{173} (2003), no.~2, 179--261.

\bibitem{Sch90}
M.~Scharlemann,
\emph{Producing reducible 3-manifolds by surgery on a knot},
Topology \textbf{29} (1990), no.~4, 481--500.

\end{thebibliography}

\end{document}
```